\scanpage{113}
% L1–8
Soient $X,Y$ deux préschémas. \uline{Application géométrique} de $E\subset X$ en $Y$ : morphisme de $\mathrm{Dis}\,(E;X)=\coprod_{x\in E}\operatorname{Spec}k(x)$ en $\mathrm{Dis}\,(Y)$. \inserted{Supposons $E$ loc. constructible.} Une application géométrique est dite \uline{constructible} si tt $x\in X$ a un voisinage ouvert \inserted{affine} $U$, et une partition de $U\cap E$ en parties loc. fermées $Z_i$, telle que $f\mid(Z_i)_{\mathrm{réd}}$ soit un \uline{morphisme} $Z_i\longrightarrow Y$. \deleted{Il revient aussi au =
de dire que \unsure{.}}

% L8–16
Ex : une appl. géom. définie par un morphisme. Si \deleted{\unsure{on}} $E=\bigcup_{i\in I}E_i$ famille \uline{quelconque} de parties loc. constructibles, alors $f:E\xrightarrow{\text{appl. géom.}}Y$ est constr. ssi pour tt $i$, $f\mid E_i$ l'est ; cela signifie aussi, \inserted{si $\deleted{\unsure{E|Z}}X$ est quasi-cpt et qu. sép.} qu'on peut trouver une \uline{partition} globale de $E$ en parties $Z_i$ constructibles loc. fermées, telles que les $f\mid Z_i$ soient des morphismes.

% L16–18
Définition de $\operatorname{Appl}_{\mathrm{cons}}((X,E),(\deleted{\unsure{X}}Y,F))$, Transitivité. Catégorie \inserted{$(\mathrm{Cons}\ \mathrm{aff})$} des modèles constructibles affines :

% L19–24
objets les schémas affines (ou si on préfère, quasi-compacts qu. sép.) et morphismes les appl. constructibles. Si $\deleted{X}E$ est une partie loc. \deleted{\unsure{préschémas|préschéma} \unsure{loc.}} constructible d'un préschéma $X$, on lui associe un foncteur
\[
\mathrm{Cons}\,(\deleted{\unsure{E}}X,E):(\mathrm{Cons}\ \mathrm{aff})^{\circ}\longrightarrow(\mathrm{Ens})
\]
% archive sheet number: 112
